\documentclass[sn-mathphys-num,pdflatex]{sn-jnl}

\usepackage{graphicx}
\usepackage{booktabs}
\usepackage{amsmath,amssymb}

\newcommand{\Drho}{D(\rho)}
\newcommand{\rhostar}{\rho^{*}}

\begin{document}

\title[Topology of Reachable Nature]{The Topology of Reachable Nature:
Persistent Homology and the Resilience of Green-Space Access in Cities of
Contrasting Morphology}

\author[1]{\fnm{}\sur{[Author name omitted for double-blind peer review]}}

\affil[1]{[Affiliation omitted for double-blind peer review]}

\abstract{SDG~11.7 targets universal access to green space, yet accessibility is measured as a
static snapshot, silent on how it degrades when the network fails. A topological
resilience metric is introduced: the Wasserstein distance between persistent-homology
diagrams of the green-space accessibility field before and after disruption. Across
$3{,}999$ districts in eight contrasting cities, street-network morphology is
consistently associated with the resilience of the \emph{connected} green-space access
structure: longer, more circuitous streets with greater resilience, disordered
orientation with less. Effects survive a wild cluster bootstrap and spatial models and
transfer out-of-sample (leave-one-city-out mean Spearman~$0.42$), though a real
clustered flood does not reproduce them. It validates only moderately against distance-
and population-weighted benchmarks (Spearman~$0.33$); its negative relation to
connectivity loss marks it as \emph{connected-access resilience}, distinct from
percolation robustness.}

\keywords{persistent homology, topological data analysis, green-space accessibility,
urban resilience, street-network morphology, SDG 11}

\maketitle

\section{Introduction}
Target~11.7 of Sustainable Development Goal~11 commits states to provide
``universal access to safe, inclusive and accessible, green and public spaces''
\citep{un2015sdg}. Planners typically measure progress with catchment or
proximity analyses: the share of residents within a walking threshold of a park.
Such measures are static. They describe access on an ordinary day and are silent
about \emph{resilience}---how access to nature holds up when the pedestrian
network is degraded by a flood, an earthquake, or the closure of a critical
bridge or street. Two gaps follow. First, a city can score well on average
proximity yet collapse under modest disruption; average accessibility does not
reveal this fragility. Second, accessibility work rarely connects the
\emph{shape} of access to urban \emph{form}, which is precisely the
architecture--mathematics question that motivates this collection.

Persistent homology (PH) offers a principled, multi-scale, coordinate-free
language for the ``holes'' and connected pieces of a coverage pattern
\citep{otter2017roadmap,feng2021tda}. PH has been used to detect coverage gaps
for amenities such as polling sites \citep{hickok2022polling} and healthcare
\citep{gonzalez2025healthcare}, to read the topology of urban congestion
\citep{wu2017congestion}, and to quantify displacement under urban expansion
\citep{rodriguez2025displacement}. What is \emph{not} novel is using PH to find
accessibility ``holes'' for an amenity. The contribution here is to move from
\emph{detecting} coverage gaps to \emph{measuring the resilience of the access
topology} and to \emph{linking that resilience to urban form}. A scope note is in
order at the outset: what is measured is \emph{network} accessibility, the pedestrian
graph's spatial reach to green space with every network node weighted equally, and
not \emph{population} accessibility, which would weight locations by the residents
they serve. The two coincide only where population is uniform, so the claims made
here concern the network's structural capacity to keep places near nature, not how
many people retain access.

\paragraph{Contributions.}
\begin{itemize}
  \item \textbf{A topological resilience metric} (the mathematics). $\Drho$ is
    defined as the Wasserstein distance between the persistent-homology diagram of
    the green-space accessibility field before and after disruption of intensity
    $\rho$, and is summarised by an area-under-curve (AUC) and a critical
    transition $\rhostar$ (Section~\ref{sec:method}). Throughout, $\Drho$ is taken to
    measure one specific construct---the resilience of the \emph{connected} green-space
    access structure, hereafter \emph{connected-access resilience}: how stably locations
    stay organised into well-served $H_0$ components as the network is thinned, as
    distinct from total accessibility loss or global urban resilience. This term is used
    consistently for what the metric captures.
  \item \textbf{First topological treatment of SDG~11.7 green-space access} in an
    architecture--mathematics framing.
  \item \textbf{A district-level morphology$\leftrightarrow$resilience result}
    across $3{,}999$ districts in eight cities, with city fixed effects, showing a
    consistent association between street-network morphology and connected-access
    resilience whose sign is stable across disruption models and district resolutions
    (Section~\ref{sec:results}).
\end{itemize}

\section{Related work}
\emph{TDA of spatial systems.} \citet{feng2021tda} survey persistent homology for
spatial data; \citet{otter2017roadmap} give the computational foundations. Within
urban analysis, \citet{wu2017congestion} read traffic congestion as a barcode,
and a line of ``resource coverage'' work applies sublevel or network-distance
filtrations to amenity access: polling places \citep{hickok2022polling} and
sexual-health services \citep{gonzalez2025healthcare}. \citet{rodriguez2025displacement}
use PH to quantify displacement under urban expansion. These studies establish
that PH can \emph{locate} coverage holes. The object of study here differs: not
where the deserts are, but how the whole access topology \emph{degrades} under
disruption, with that degradation regressed on morphology.

\emph{Accessibility and equity.} A large body of work measures green-space access
as a static field, whether by simple footpath proximity \citep{carthy2020greenspace}
or by catchment methods such as the Gaussian two-step floating catchment area
\citep{sun2023catchment}, often to diagnose inequities in who lives close to nature
\citep{asibey2026equity}. The accessibility field used here (Section~\ref{sec:method})
is built the same way (network walking time to the nearest green space) but is then
read topologically and stress-tested, rather than thresholded once.

\emph{Street-network morphology.} Form is characterised with standard descriptors
computed from OpenStreetMap via OSMnx \citep{boeing2017osmnx}: intersection
density, circuity, orientation entropy, and mean street-segment length. Relating
these to a resilience outcome is the architecture--mathematics bridge of this
paper.

\section{Mathematical background}
Let $G=(V,E)$ be the pedestrian network and $f:V\to\mathbb{R}$ a scalar field
giving, at each node, the network walking time to the nearest green-space access
point. The \emph{sublevel-set filtration} of $f$ builds a growing simplicial
complex $K_t = \{\,\sigma : \max_{v\in\sigma} f(v)\le t\,\}$. As the threshold
$t$ (walk-minutes) rises, well-served regions appear and merge. Persistent
homology records, for each topological feature, a birth $b$ and death $d$; the
multiset of points $(b,d)$ is the \emph{persistence diagram} $\mathrm{Dgm}_k(f)$
in homological dimension $k$ \citep{edelsbrunner2000persistence}. Dimension~0
tracks connected components (the merging of well-served regions); dimension~1
tracks loops (enclosed voids, i.e.\ access ``deserts'').

The distance between two diagrams is measured by an optimal matching. The
$1$-Wasserstein distance is
\[
W_1(X,Y)=\min_{\gamma}\sum_{x\in X}\lVert x-\gamma(x)\rVert,
\]
minimised over bijections $\gamma$ between the diagrams (with matching to the
diagonal allowed). For a \emph{fixed} domain, these diagram distances are stable to
perturbations of the field $f$: the bottleneck distance by the classical stability
theorem \citep{cohensteiner2007stability}, and the $1$-Wasserstein distance used here
by the $L_p$-stability result of \citet{cohensteiner2010lpstable}. A caveat must be
stated precisely here, because it bounds what stability can justify. Both theorems
bound the diagram distance by $\lVert f-g\rVert_\infty$ \emph{for two functions on the
same simplicial complex} $K$: they control the effect of perturbing the \emph{field}
while the domain is held fixed. Edge removal is not of this form. Removing edges
replaces the complex $K$ by a proper subcomplex $K'\subsetneq K$ and only then
recomputes the walk-time field on $K'$, so the baseline and disrupted diagrams are
sublevel-persistence of \emph{different} filtrations, $f$ on $K$ and $f'$ on $K'$, and
there is no single ambient function whose sup-norm the stability inequality could
apply to. Forcing a common domain---extending $f'$ to all of $K$ by assigning the
removed structure, and any node thereby cut off from green space, the value
$+\infty$---makes the mismatch explicit: points of the diagram escape to infinity and
leave the finite $1$-Wasserstein matching, so the finite diagram distance is not
bounded by any $\lVert f-f'\rVert$ term and can even \emph{decrease} as mass departs
(the mechanism behind the weak/negative connectivity relation in
Section~\ref{sec:results}). The classical stability theorem therefore does
\emph{not} directly bound $\Drho$: it rigorously controls only the one part of the
disruption that is a pure field change on a fixed domain---re-snapping the
accessibility field where the surviving subgraph coincides with the original---while
the genuine domain change $K\to K'$ carries no such guarantee, so the whole map from
disruption intensity to diagram distance need not be well-behaved. In particular, stability does not imply
monotonicity, and $\Drho$ is indeed \emph{not} always monotone: at high $\rho$ the
İstanbul and Phoenix curves turn down as the network fragments and finite classes
vanish (Section~\ref{sec:results}). Stability is thus a mathematical well-posedness
property of the field-to-diagram map, not a guarantee about the metric under network
damage; the metric is treated throughout as a \emph{proposed proxy} for how the
access topology holds up, to be validated against independent resilience evidence in
future work rather than as ground truth. PH is computed with GUDHI \citep{gudhi} and diagram distances with
\texttt{persim}.

\section{Methodology}\label{sec:method}
\paragraph{Data.} For each city the pedestrian network and green spaces are taken
from OpenStreetMap via OSMnx \citep{boeing2017osmnx}, and the
analysis extent is fixed to the city's \emph{urban centre} polygon from the GHS Urban
Centre Database \citep{ghsl2019ucdb} so that extents are comparable across
countries. Green spaces are the OSM classes \texttt{leisure=park|garden|recreation\_ground},
\texttt{landuse=grass|recreation\_ground}, \texttt{place=square} and
\texttt{natural=wood}; their centroids, snapped to the nearest network node, are
the access points. These are termed \emph{OSM-mapped green} spaces, not \emph{public}
spaces: OSM tags do not certify public accessibility, so a private garden, a fenced
lawn or an impenetrable wood can enter the set. The classes are included because
together they capture the everyday green and open spaces a pedestrian can plausibly
reach; the trade-off is heterogeneity, addressed by the class-exclusion sensitivity
below. This is a deliberately inclusive definition: it
merges spaces of very different type, size and usability (a pocket garden, a city
square and a large wood all count equally), and it takes the centroid as a single
access point rather than modelling entrances or area. A sensitivity analysis under a
strict, accessibility-conservative subset (dropping wood, grass and gardens; keeping
parks, recreation grounds and squares) is reported below. Every green space is treated
as an equivalent point source; distinguishing quality, size, or public/private status
is left to future work (see Limitations).

\paragraph{Accessibility field.} Each node is assigned its network walking time to the
nearest green-space access point by multi-source Dijkstra, giving the field $f$
(Fig.~\ref{fig:accessfield}). Length is converted to time at a constant
$4.8$~km/h ($1.33$~m/s), which sits within the comfortable adult gait speed of
$1.27$--$1.46$~m/s reported by \citet{bohannon1997walking}. The exact value does not
matter for the conclusions: speed enters only through $\text{time}=\text{length}/
\text{speed}$, so a different speed rescales the whole field, and hence every
diagram, Wasserstein distance and AUC, by the same constant. The regression signs,
significances and the city ranking are therefore invariant to the speed choice (if
the persistence threshold is rescaled with speed the invariance is exact). This is
confirmed empirically: on Amsterdam, recomputing at $4.0$ and $5.0$~km/h leaves
the district AUC a near-exact rescaling of the $4.8$~km/h values (Spearman
$\rho=0.98$ and $0.996$; median AUC ratios $1.24$ and $0.96$, matching the expected
$4.8/4.0$ and $4.8/5.0$). The field is defined over \emph{network nodes}, each weighted equally;
it is a property of the pedestrian network's spatial reach to green space, not a
population-weighted or distributional measure. The metric therefore speaks to one
facet of SDG~11.7, the network's capacity to keep locations close to green space and
how that capacity degrades under disruption, and not to the target's
population-coverage or inclusivity dimensions, which would require census and
land-use data not used here.

\begin{figure}[t]\centering
\fbox{\parbox[c][4.5cm][c]{0.72\linewidth}{\centering\ttfamily fig\_accessfield\\[2pt]\normalfont\itshape (drag-drop this PDF into figures/, then restore \textbackslash includegraphics)}}
\caption{Green-space accessibility field for Amsterdam: each of the $49{,}202$
pedestrian-network nodes inside the GHS-UCDB urban centre, coloured by network
walking time to the nearest green space. Green marks well-served locations,
red the access-poor pockets whose merging structure the resilience metric tracks.}
\label{fig:accessfield}
\end{figure}

\paragraph{Resilience metric.} Let $\mathrm{Dgm}(f)$ be the baseline diagram. A
disruption of intensity $\rho\in[0,1]$ removes a fraction $\rho$ of network
edges; the field and diagram are recomputed on the disrupted network, defining the
degradation
\[
\Drho=\underset{r}{\mathbb{E}}\;W_1\!\big(\mathrm{Dgm}(f),\,\mathrm{Dgm}(f_{\rho,r})\big),
\]
averaged over $r$ replicate disruptions. Three disruption families are considered:
uniformly random edge removal (percolation-style), targeted removal of
high-betweenness edges, and hazard removal of edges in low-elevation zones. These
are deliberately \emph{stylised stress tests}, not simulations of specific
disasters: random removal probes generic structural robustness (and, being spatially
uniform, does \emph{not} reproduce the spatially clustered damage of a real flood or
earthquake), targeted removal probes the worst case of losing critical links, and
the hazard proxy adds a first, mild topographic signal. They serve to compare
places under a common, reproducible perturbation, and the results are read
correspondingly cautiously as real-disaster resilience (Limitations). Exact
edge betweenness is $O(V\lvert E\rvert)$ and intractable at city scale, so the
targeted scenario ranks edges by betweenness approximated from $k=500$ sampled
pivot sources. The
resilience of a place is summarised by the area under the curve
$\mathrm{AUC}=\int_0^{\rho_{\max}}\Drho\,d\rho$ (normalised by the $\rho$-range)
and the critical transition $\rhostar$, the smallest $\rho$ at which $\Drho$ first
reaches half its maximum. The half-maximum serves as a scale-free, monotone
summary of \emph{where} the degradation curve turns over: unlike a fixed absolute
threshold it adapts to each place's own range, and unlike the inflection point it
is robust to the curve's noise. It is a descriptive locator of the transition, not
a claim of a sharp percolation threshold, and (like AUC) it depends on the chosen
disruption range $[0,\rho_{\max}]$; $\rhostar$ is therefore compared only within a
fixed range.

\paragraph{Homology dimension and denoising.} On real street networks the
sublevel filtration's finite $H_1$ is nearly empty: a street graph is close to
planar and contains few triangles, so the network's cycles are never filled and
almost all $H_1$ classes are essential (infinite) and carry no finite persistence.
In Amsterdam, for example, the baseline diagram has $6{,}449$ finite $H_0$ classes
but only a single finite $H_1$ class (Fig.~\ref{fig:persistence}). This is a joint
property of the \emph{representation} (a node-valued walk-time field on the
near-planar pedestrian graph) and the \emph{sublevel filtration}, not a universal
fact about green-space access. A different construction, for instance a
Vietoris--Rips filtration on straight-line distances or a $2$-complex that fills
faces, would populate $H_1$ and could shift the signal into loops. Under this
construction the finite information lives in $H_0$ (the merging of well-served
regions), so $\Drho$ is computed on $H_0$ and the $H_1$ scarcity is treated as a
modelling consequence to revisit with alternative filtrations, not as evidence that
access ``deserts'' are absent. To keep the exact Wasserstein matching tractable
on city-scale diagrams, $H_0$ classes with lifetime below $1$~walk-minute are discarded.
This cut is principled rather than arbitrary: a one-minute difference in walk time
is below the resolution at which the constant-speed ($4.8$~km/h) field is
meaningful, so sub-minute classes are measurement noise, not access structure. Its
effect is small and was tested directly (Persistence-threshold sensitivity, below):
varying the cut over $0.5$--$2.0$ minutes preserves every sign and the significance
of the three headline descriptors, changing only reported magnitudes.

\begin{figure}[t]\centering
\fbox{\parbox[c][4.5cm][c]{0.72\linewidth}{\centering\ttfamily fig\_persistence\\[2pt]\normalfont\itshape (drag-drop this PDF into figures/, then restore \textbackslash includegraphics)}}
\caption{Persistence diagram of the Amsterdam baseline accessibility field. The
finite part is carried almost entirely by $H_0$ ($6{,}449$ classes---the merging
of well-served regions as the walk-time threshold rises) with only a single finite
$H_1$ class, motivating an $H_0$-based resilience metric.}
\label{fig:persistence}
\end{figure}

\paragraph{Districts and inference.} Statistical inference rests not on eight
cities but on intra-urban \emph{districts}: each urban-centre boundary is tiled
with a uniform H3 hexagonal grid (resolution~8, $\approx0.7$~km$^2$ cells) and,
per district, its local resilience summaries and its local morphology are computed
(intersection density, circuity, orientation entropy, mean segment length, and a
green-space fragmentation index---the number of distinct green-space patches per
square kilometre of green area within the district, computed on the green space
clipped to each hex so that it varies between districts rather than being a
per-city constant collinear with the fixed effects). The following model is then fit
\[
\mathrm{AUC}_i = \beta^\top \mathbf{m}_i + \alpha_{c(i)} + \varepsilon_i,
\]
an ordinary-least-squares regression of district AUC on morphology
$\mathbf{m}_i$ with \emph{city fixed effects} $\alpha_{c(i)}$ absorbing all
\emph{additive between-city} differences (any confound constant within a city, such
as climate, planning regime or data vintage). They do not remove within-city
confounding or effects that vary across districts inside a city; those are addressed
separately with explicit controls and spatial models below. The cities serve
as a typological overlay, not as the inferential sample. This design also sidesteps the cross-city comparability
problem of AUC. Because $\alpha_{c(i)}$ absorbs each city's overall AUC level, which
scales with network and diagram size, the morphology coefficients are identified
purely from \emph{within-city} variation, and the inference never rests on comparing
absolute AUC across cities. Only the descriptive city typology, not the regression,
would need a size normalisation. For inference that respects the nesting, standard
errors are also clustered at the \emph{city} level---the eight clusters are the eight
cities, the level at which each disruption and each fixed effect operate, and within
which districts share a network, climate and planning regime. Eight clusters is a
small number, at which the asymptotic cluster-robust $p$-values are only approximate;
they are therefore reported alongside a wild cluster bootstrap
(Appendix~\ref{app:robust}), and the cities are treated as the practical, if not
strictly independent, units of between-city variation.

\section{Case studies}
Eight cities on four continents are studied, chosen to span street morphology, planning
regime and the Global North/South divide: İstanbul (organic, hilly, polycentric,
seismic), Barcelona (Cerdà grid with superblock interventions), Amsterdam (compact,
canal, walkable), Bogotá (dense informal--formal mix), Phoenix (low-density grid
sprawl), Singapore (planned, tropical), Nairobi (rapidly growing, mixed-formality)
and Vienna (compact, central European on the Danube). This is a \emph{purposive
maximum-contrast} sample, selected to maximise morphological and planning diversity so
that any consistent morphology--resilience relationship is stress-tested against very
different urban forms; it is explicitly \emph{not} a probability sample of world
cities, so the findings are framed as holding for these eight cities and are tested for
out-of-sample transfer (below) rather than asserted as universal. Boundaries are the GHS-UCDB
urban-centre polygons; homonymous centres (e.g.\ Barcelona, Spain vs.\ Venezuela)
are disambiguated by largest true (equal-area) extent. Two further candidates
(Copenhagen, Melbourne) were dropped because the OpenStreetMap place query did not
resolve their full urban extent. Some robustness and validation analyses
(Section~\ref{sec:results}, Appendix~\ref{app:robust}) were computed on the original
five-city core and are labelled as such.

\section{Results}\label{sec:results}
The random-disruption analysis yields $3{,}999$ qualifying districts (147--962 per
city, Table~\ref{tab:typology}); only $2.6\%$ have zero AUC. District mean AUC ranks
Amsterdam as the most resilient and İstanbul/Bogotá as the least; Nairobi has the
earliest breakdown ($\rhostar=0.11$). The district-mean and whole-city rankings
agree at the extremes (Amsterdam most resilient, Bogotá least) but not in the
middle---Barcelona is mid-ranked by district mean yet most resilient at the
whole-city scale---a reminder that AUC is comparable across cities only in broad
ordering, not level.

\begin{table}[t]\centering
\caption{District coverage and city typology (descriptive; H3 resolution~8,
$n=3{,}999$ across eight cities). Higher mean AUC $\Rightarrow$ access topology
changes more under disruption $\Rightarrow$ \emph{less} resilient.}
\label{tab:typology}
\begin{tabular}{lcccc}
\toprule
City & districts & mean AUC & mean $\rhostar$ & mean $H_0$ tot.\ persist.\\
\midrule
Amsterdam & 227 & 4.46 & 0.25 & 13.11\\
Nairobi   & 363 & 5.90 & 0.11 & 9.09\\
Vienna    & 460 & 6.08 & 0.22 & 12.91\\
Phoenix   & 962 & 6.68 & 0.16 & 8.08\\
Singapore & 573 & 7.34 & 0.17 & 11.96\\
Barcelona & 147 & 7.67 & 0.28 & 13.63\\
İstanbul  & 729 & 7.75 & 0.20 & 12.84\\
Bogotá    & 538 & 8.02 & 0.24 & 16.90\\
\bottomrule
\end{tabular}
\end{table}

The headline result is the district-level regression (Table~\ref{tab:regression}),
fit over eight cities and the full disruption grid ($8$ intensities $\times\,10$
replicates). With city fixed effects and $n=3{,}999$, three of five morphology
descriptors are significant, all at $p<10^{-16}$. Because the coefficient sign is on
AUC (higher AUC = less resilient), a negative coefficient means the descriptor is
associated with \emph{greater} resilience. More circuitous networks (coef $-7.39$)
and longer street segments ($-0.050$) are associated with more resilient green-space
access, whereas higher orientation entropy ($+1.85$) is associated with less
(Fig.~\ref{fig:morphology}). Because the raw coefficients sit on different scales, magnitude is
read from the standardised $\beta$ and the $10$th--$90$th percentile effect
(Table~\ref{tab:regression}): \emph{mean street length} is the largest effect (std.\
$\beta=-0.28$; a $10$th-to-$90$th-percentile change moves AUC by $62\%$ of its
standard deviation), then orientation entropy ($+0.16$) and circuity ($-0.13$); the
tiny $p$-values reflect the large sample, not effect magnitude. The descriptors are
not collinear (all variance-inflation factors $<3$), and the pooled signs are not a
pooling artefact---each of the three holds in seven of the eight cities. Intersection
density is not significant ($p=0.71$) and its sign is unstable across samples, and
green-space fragmentation is likewise not significant. Districts nest within the eight
cities and are spatially contiguous, so their residuals are not independent;
Appendix~\ref{app:robust} treats this in full. Under a wild cluster bootstrap---the
appropriate small-cluster correction for the eight clusters (the eight cities)---
circuity and mean street length remain significant ($p=0.0002$ each), while
orientation entropy does not ($p=0.16$); orientation entropy is nonetheless strong in
the spatial-error and spatial-lag models ($|z|\approx10$). Circuity
and mean street length are therefore treated as the robust pair and orientation entropy as a real but
less certain third effect.

\begin{table}[t]\centering
\caption{District fixed-effects regression, $\mathrm{AUC}\sim\text{morphology}+C(\text{city})$,
$n=3{,}999$ across eight cities, random-disruption scenario at the full disruption
grid ($8\times10$). The standardised coefficient (std.\ $\beta$) and the
change in AUC across each descriptor's $10$th--$90$th percentile (as \% of the AUC
standard deviation) so that effects on different scales are comparable and not read
off the raw coefficient or the (large-$n$) $p$-value alone. See Table~\ref{tab:spatial}
for city-clustered and spatial-model inference. Sig.: $^{***}p<0.001$.}
\label{tab:regression}
\begin{tabular}{lrrrr}
\toprule
Morphology descriptor & coef & std.\ $\beta$ & $\Delta_{10\text{--}90}$ (\% SD) & $p$-value\\
\midrule
mean street length & $-0.0505$ & $-0.28$ & $-62$ & $4.9\times10^{-29}$ $^{***}$\\
orientation entropy & $+1.845$ & $+0.16$ & $+51$ & $1.4\times10^{-27}$ $^{***}$\\
circuity & $-7.387$ & $-0.13$ & $-27$ & $3.6\times10^{-17}$ $^{***}$\\
intersection density & $+0.021$ & $+0.01$ & $+2$ & $0.71$\\
green-space fragmentation & $-3.7\times10^{-7}$ & $-0.00$ & $-0$ & $0.79$\\
\bottomrule
\end{tabular}
\end{table}

\begin{figure}[t]\centering
\fbox{\parbox[c][4.5cm][c]{0.72\linewidth}{\centering\ttfamily fig6\_morphology\\[2pt]\normalfont\itshape (drag-drop this PDF into figures/, then restore \textbackslash includegraphics)}}
\caption{Headline morphology$\leftrightarrow$resilience relationship. Each point
is one of the $3{,}999$ districts across eight cities, coloured by city; the dashed
line is the pooled least-squares fit. Higher AUC indicates a district whose
green-space access topology degrades more under disruption (less resilient).}
\label{fig:morphology}
\end{figure}

The whole-city scale gives an \emph{illustrative} typology, not a second
inferential result: it carries no statistical weight, and it is reported to visualise
the district-level signal, not to rank cities definitively. For the five original
cities, computed on the same UCDB-clipped networks, the degradation curves $\Drho$
rise smoothly and order the cities consistently with the district means
(Fig.~\ref{fig:resilience}). Because city-level AUC is computed on the full-city
diagram (thousands of $H_0$ classes) it is two to three orders of magnitude larger
than the per-district means of Table~\ref{tab:typology} and scales with network
size; it is \emph{not} normalised, so only the relative ordering---never the
level---should be read, and a size-normalised cross-city AUC is left to future work.
On that ordering, city-level AUC is
lowest, and so most resilient, for Barcelona ($620$) and Amsterdam ($898$), and highest
for İstanbul ($2440$) and Bogotá ($2828$), with Phoenix in between ($2053$). The
cities that would informally be called compact, Barcelona and Amsterdam, retain more of
their access topology under disruption than the large, sprawling or polycentric
ones; this pattern is described with the measured morphology descriptors rather
than with ``walkability,'' which is not modelled directly. Under
extreme removal ($\rho=0.5$) the İstanbul and Phoenix curves turn down as the
network fragments so heavily that unreachable nodes drop out of the finite
diagram, a signature of near-total collapse rather than of recovered access.

\begin{figure}[t]\centering
\fbox{\parbox[c][4.5cm][c]{0.72\linewidth}{\centering\ttfamily fig5\_resilience\\[2pt]\normalfont\itshape (drag-drop this PDF into figures/, then restore \textbackslash includegraphics)}}
\caption{City-level resilience curves $\Drho$ for the five original cities
(random-disruption scenario), computed on the UCDB-clipped networks; illustrative.
Dotted vertical lines mark each city's critical transition $\rhostar$.}
\label{fig:resilience}
\end{figure}

\paragraph{Robustness (summary).} The morphology relationship is stable across every
alternative specification tried; Appendix~\ref{app:robust} gives the full results
(the additional sensitivity checks were run on the original five-city core and are
labelled there) and Fig.~\ref{fig:robustness} summarises them. Briefly: the three
headline signs and
their significance ($p<10^{-3}$ or better) survive a \emph{targeted} betweenness
disruption, the coarse $5\times3$ and full
$8\times10$ disruption grids, H3 resolutions $7$/$8$/$9$ (a $28\times$ range of
district counts), persistence thresholds from $0.5$ to $2.0$~walk-minutes, a halving
and doubling of the betweenness pivot count, and the inclusion of four
district-level controls (green-space area, road hierarchy, relief and city-centre
distance). Intersection density is the one descriptor that moves: significant under
some specifications but not the full-grid headline, so it is treated as the weakest,
grid-sensitive effect throughout. The one disruption family that does \emph{not}
carry the signal is the spatially clustered real-flood hazard, which is treated as a
distinct question (below).

\paragraph{Green-space definition sensitivity (five-city core).} Because the inclusive green-space set
mixes reliably public parks with possibly-private gardens, fenced grass and
impenetrable woods, the analysis was re-run on a strict, accessibility-conservative
subset: only \texttt{leisure=park}, \texttt{recreation\_ground} and
\texttt{place=square}, dropping wood, grass and gardens. This is a large cut (in
Amsterdam, from $23{,}657$ features to $369$), yet the district AUC is largely
preserved (Spearman $0.82$ against the inclusive set on the same reduced grid), and
the regression is unchanged in substance: circuity ($-8.08$, $p=4\times10^{-9}$),
mean street length ($-0.048$, $p=2\times10^{-12}$) and orientation entropy ($+1.21$,
$p=2\times10^{-7}$) keep their sign and significance, and intersection density
remains the weak, non-robust term. The morphology result is thus not an artefact of
counting inaccessible green as accessible.

\paragraph{Criterion validity: benchmark against an accessibility measure (all eight cities).} If the
topological machinery is to be worth its cost it should track a simpler measure of
access loss. It is benchmarked, per district, against a \emph{distance-based}
accessibility-degradation measure, namely the increase in mean walk-time to green
space under the same disruption (a proximity-accessibility degradation of the kind
used throughout the field). As a second reference, a \emph{connectivity}
degradation, the increase in the fraction of nodes left unable to reach any green
space; this is the natural non-persistent counterpart, since $H_0$ of a sublevel
filtration is single-linkage connectivity. The level of comparison is what settles
the question. The inference identifies morphology from \emph{within-city} variation
(city fixed effects), so the matching validity question is also within-city, once
each city's overall AUC level (which scales with network size) is removed. Within
city, and the benchmark is computed for all \emph{eight} headline cities, so the
validation sample matches the analysis sample. The topological AUC correlates
positively with the distance benchmark, with a \emph{moderate} pooled magnitude
(within-city pooled Spearman $\rho=+0.33$ over the $3{,}999$ districts; the tiny
$p<10^{-102}$ reflects the sample size, not strength). The relationship is, however,
markedly \emph{heterogeneous}: positive in seven of the eight cities but strong in
only one---Amsterdam $+0.47$---modest in Bogotá $+0.16$ and İstanbul $+0.10$, near
zero in Nairobi $+0.09$, Vienna $+0.07$, Barcelona $+0.05$ and Phoenix $+0.03$, and
marginally negative in Singapore $-0.03$ (not significant). This is moderate
supporting evidence---not strong validation---that the metric tracks a simple
accessibility-degradation ground truth \emph{on average} at the level at which it is
used, while being an unreliable graded-accessibility proxy within any single city. The naive pooled correlation is much
weaker ($\rho=+0.07$) only because cross-city AUC scale swamps it, a Simpson's
paradox driven by the same scale the fixed effects absorb. The benchmark also holds
when access is \emph{population-weighted} rather than node-uniform: weighting each
node by the open GHS-POP residential-population grid \citep{ghsl2023pop} (JRC,
$1$~km, $2020$; the same GHSL family and Mollweide grid as the urban-centre
boundaries), computed for all eight cities, gives an almost identical within-city
correlation ($\rho=+0.33$, $p<10^{-101}$). So the topological signal is not an artefact of weighting locations
equally, and it speaks, at least for graded access, to the population-coverage
dimension of SDG~11.7 and not only to bare network reach. Two qualifications
apply. The metric adds value rather than echoing the benchmark: controlling for
both baselines and city fixed effects, the three headline morphology descriptors
stay significant ($p<10^{-6}$), so it carries configurational information the simple
measures miss. And the convergent validity is with \emph{graded} access loss, not
with \emph{disconnection}: against the connectivity-degradation reference the
eight-city correlation is significantly \emph{negative} ($\rho=-0.21$,
$p\approx10^{-42}$), because disconnected nodes leave the finite diagram. Rather than
a failure, this negative relation is \emph{discriminant} validity---it establishes
that $\Drho$ is not a repackaged connectivity or percolation-robustness measure but
captures a distinct construct. Precisely, $\Drho$ measures the
\emph{graded topological persistence of the connected green-access structure under
progressive network thinning}: how stably locations stay organised into well-served,
mutually close $H_0$ components as edges are removed, rather than whether they stay
connected at all. It is thus validly read as the resilience of the \emph{connected}
served structure and should be paired with a disconnection count for total access
loss; the Limitations show that neither a naive finite cap nor extended persistence
folds the two into one scalar cleanly (both recover the connectivity sign only by
sacrificing the graded signal), so the two channels are kept separate and the claim
is scoped accordingly. Reporting the disconnection channel as an explicit
\emph{companion outcome} makes the complementarity concrete: regressing the
connectivity-loss AUC on the same morphology and fixed effects, street form predicts
outright disconnection too ($R^2=0.44$) but with a different and sometimes
\emph{opposing} signature: more circuitous districts are graded-more-
resilient yet also \emph{more} disconnection-prone (standardised $\beta=+0.15$,
$p<10^{-28}$ on the connectivity channel). The two are not substitutes: a full account
of access resilience must report both the graded-persistence and the connectivity-loss
channels, as done here, and the metric's ``resilience'' is exactly the former,
\emph{connected-access resilience}, not global accessibility or urban resilience.

\paragraph{Mechanism: is the association mediated by path redundancy?} The natural
reading of the circuity and segment-length effects is that they proxy
\emph{alternative-path redundancy}---more ways to reach the same green space, so that
losing one link rarely isolates a served region. This mechanism is testable directly.
For each district a structural redundancy measure is computed from the network alone,
independent of the disruption simulation and the accessibility field (so it is not
circular with $\Drho$): the \emph{meshedness} (the $\alpha$ index, the density of
independent cycles, $\alpha=(e-v+p)/(2v-5)$, the standard planar-network route-
availability measure \citep{barthelemy2011spatial}), and, as a coarse check, mean
node degree. The direct test does \emph{not} support the redundancy story. Within
city, redundancy correlates with \emph{lower} resilience, not higher: meshedness
against the degradation AUC gives Spearman $\rho=+0.50$ and mean degree $\rho=+0.60$
(both $p<10^{-250}$; recall higher AUC is less resilient). Morphology predicts
redundancy with the sign that \emph{opposes} the mechanism---more circuitous
($-0.045$, $p=10^{-4}$) and longer-segment ($-6\times10^{-4}$, $p<10^{-19}$) districts
are \emph{less} mesh-redundant---and adding meshedness to the headline regression
leaves the morphology effects essentially intact (circuity attenuates by $10\%$,
segment length by $18\%$, orientation entropy not at all), while meshedness itself
enters with the counter-mechanism sign ($+16.4$, $p<10^{-44}$; $R^2$ $0.196\to0.235$).
The metric therefore rewards \emph{coarse}, low-cycle-density configurations (long,
circuitous segments, few redundant loops) whose served $H_0$ structure reorganises
little under random thinning, rather than finely meshed grids with many alternative
paths; part of this may be that denser districts carry more $H_0$ classes and hence
larger Wasserstein mass. Either way, the simple alternative-path account is refuted,
and the mechanistic language below is kept at the level of interpretation.

\section{Discussion}
The results give an interpretable, architecture-facing reading of connected-access
resilience.
The robust effects are circuity and segment length: districts with more circuitous
routing and longer segments retain more of their well-served structure when edges
are removed. The intuitive explanation---that alternative paths to green space survive
the loss of any one link---is, however, \emph{not} borne out by the direct redundancy
test above: a standard mesh-redundancy index runs the wrong way, so the effect is
better read as coarse, low-cycle configurations producing a more stable $H_0$
signature than as path substitution. Intersection density is not significant, and its
sign is not even stable
between the five- and eight-city samples, so no conclusion is drawn from it. Districts
with more disordered street orientation (higher entropy) are more brittle. That the signal
is carried by $H_0$ rather than $H_1$ has a concrete planning meaning. In
green-dense cities the structure that matters is not isolated ``deserts''
surrounded by served land but the connectivity of the served regions themselves,
and that connectivity is what disruption erodes.

For design, this suggests a \emph{hypothesis} rather than a prescription, and one the
present data qualify. One reading---that \emph{path redundancy}, not raw
proximity, protects access to nature under stress---was tested directly and not
supported: the descriptors associated with resilience (circuitous routing, longer
segments) are in fact associated with \emph{less} mesh redundancy, not more. What the
metric rewards is a coarser, more stable served configuration under random link loss.
Whether that stability is a genuine planning virtue or partly an artefact of how the
topological summary scales with configuration granularity is what an
independent, validated resilience outcome would have to settle. Any reframing of an
SDG~11.7 design brief from ``put a park within $x$ metres'' toward ``ensure the walk
to it survives the loss of individual streets'' therefore remains a direction for
testing: the association is observational, the outcome is a proposed topological
indicator rather than a validated resilience measure, its own proposed mechanism does
not survive a direct test, and green-space quality, public accessibility and real
hazard exposure are not modelled.

\paragraph{Limitations.} The caveats are grouped by the kind of claim they qualify.

\emph{Inference and spatial dependence.} The design is cross-sectional and
observational, so the regression identifies \emph{associations} between morphology
and topological resilience, not causal effects. Districts are also not independent:
they nest within eight cities and are spatially contiguous. City fixed effects
absorb additive between-city differences, but two residual problems remain. First,
eight cities furnish only eight clusters for cross-city inference; clustering the
standard errors by city widens the intervals, but all three headline effects still
remain significant under this conservative treatment (circuity $p=2\times10^{-4}$,
mean street length $p=5\times10^{-19}$, orientation entropy $p=0.036$;
Table~\ref{tab:spatial}). Eight clusters is nonetheless modest, so greater weight is placed on the
estimated spatial models rather than on small-cluster asymptotics. Second, the
residuals retain weak positive spatial autocorrelation within most cities (global
Moran's~$I$ up to $0.18$; pooled $0.24$, $p<10^{-100}$). This is addressed by
estimating spatial-error and spatial-lag models on a within-city weights matrix
(Table~\ref{tab:spatial}): both absorb strong spatial dependence
($\hat\lambda=0.53$, $\hat\rho=0.51$), yet all three headline descriptors keep their
sign and remain significant at $|z|\ge5$. The morphology signal is therefore not an
artefact of spatial autocorrelation. Broader replication---more cities, and cities
with sparser OpenStreetMap coverage than this set---remains the natural next step.

\emph{Omitted variables.} Morphology is only one of many district attributes that
could be associated with resilience. The four that the data
support are controlled---green-space
area, road hierarchy, relief and city-centre distance (Table~\ref{tab:controls})---
and the headline morphology effects survive their inclusion, so the association is
not explained away by these four (which is weaker than ruling them out as causes).
But population and built density, land-use mix, park-level size
and quality, socioeconomic status and pedestrian-infrastructure quality would
require external data (e.g.\ census and GHSL built-up layers) and remain
uncontrolled; city fixed effects remove their between-city component but not their
within-city variation, so some residual omitted-variable bias may persist.

\emph{Shared origin of predictor and outcome.} Both the morphology descriptors and
the resilience outcome are computed from the \emph{same} street network, so part of
the association is structural rather than a clean independent-variable/dependent-
variable relationship: a district cannot have, say, low circuity and a routing
topology that behaves as if it were highly circuitous. The
regression is therefore read as characterising \emph{which} formal properties co-vary with resilient
access, not as an external cause acting on an independent outcome. Two observations
temper the concern: the morphology descriptors explain only a modest share of the
variance ($R^2\approx0.20$), so the outcome is far from a
deterministic re-encoding of the predictors; and the accessibility field also
depends on the green-space configuration, which is exogenous to street form. A direct
check is available. Because the AUC scales strongly with the raw size of the
persistence diagram (its correlation with the district's total $H_0$ persistence is
$+0.83$), the morphology signal is refitted while controlling for that
same-network measure of diagram magnitude. Circuity and mean street length survive
($-7.39\to-2.88$ and $-0.050\to-0.014$, both still $p<10^{-8}$), so they carry
configurational information beyond the sheer size of the diagram; orientation entropy,
by contrast, loses significance and flips sign ($+1.85\to-0.14$, $p=0.10$), revealing
that its apparent effect was largely a re-encoding of diagram magnitude---a concrete
instance of the shared-origin risk. This is only a partial mitigation: the control is
itself network-derived, and a design that predicted one city's resilience from
another's morphology, or from an independent disruption dataset, would be
needed to break the shared-origin coupling entirely.

\emph{What the metric measures, and how far it is validated.} $\Drho$ is a
\emph{proposed proxy} for connected-access resilience, and this is the study's central
limitation. What is established rigorously is an \emph{association} between morphology
and this particular Wasserstein/AUC topological indicator; that circuity and mean
street length ``explain resilience'' in a broader sense follows only to the extent
that the indicator is itself a valid measure of resilience, which is not fully
demonstrated. Evidence of two kinds is offered. Its
\emph{construct validity} is supported by its behaviour---it rises with disruption
over most of the range (though not monotonically at high $\rho$, where finite classes
vanish), builds on field-perturbation stability of the diagram distance on the fixed
domain \citep{cohensteiner2007stability} (with the domain-change caveat of
Section~\ref{sec:method}), and yields a morphology relationship that survives
disruption model, district resolution, denoising threshold, pivot count,
spatial-error/lag modelling and four omitted-variable controls. On \emph{criterion} validity, the benchmark against a
distance-based accessibility-degradation measure, computed for all eight headline
cities so that validation and analysis samples coincide, is positive at the level at
which the metric is used but only moderately and heterogeneously so. Within city, the
unit the fixed effects identify from, the topological AUC correlates $+0.33$
($p<10^{-102}$) with the benchmark pooled over the $3{,}999$ districts and is positive
in seven of the eight cities, yet strong in only one (Amsterdam $+0.47$) and near zero
in most; it remains non-redundant (the morphology signal survives controlling for it).
Its \emph{discriminant} validity is cleaner: against connectivity degradation the
eight-city correlation is significantly negative ($\rho=-0.21$), which shows the
metric is not a restatement of connectivity or percolation robustness but a distinct
construct---the graded topological persistence of the \emph{connected} green-access
structure under progressive network thinning. That construct has a definite scope: it
captures graded access degradation, not total \emph{disconnection} (unreachable nodes
leave the finite diagram), so the metric should be paired with a disconnection count
for total access loss. An attempt to fold disconnection into a
single scalar is reported, and the result is informative. A finite cap ($\tau=60$~min) recovers the
connectivity sign ($-0.10\to+0.02$ on a two-city check) but erodes the
graded-distance correlation ($+0.41\to+0.06$), and a principled extended-persistence
diagram, which gives the disconnected components finite coordinates in a separate
channel, behaves the same way on Amsterdam (connectivity $+0.31$ but distance only
$+0.05$), because the high-persistence component classes dominate a single
Wasserstein match. Graded resilience and disconnection are, then, two channels that
a single scalar conflates; reporting them separately, as done here, is the right
treatment. A joint multi-channel metric, and a full head-to-head against an
externally \emph{validated} criterion (a formal two-step floating catchment
accessibility model, or observed post-disruption service loss) rather than the
convergent proxies benchmarked here, is the natural next step. Two further
caveats:
AUC depends on network and district size and on the number of $H_0$ classes, so it
is comparable across cities only in relative ranking, not in absolute level (city
fixed effects absorb this scale in the regression), and $\rhostar$ depends on the
chosen disruption range. The down-turn of the İstanbul and Phoenix curves at
$\rho=0.5$ reflects finite nodes dropping out of the diagram under near-total
fragmentation, not recovered access, and marks the upper edge of the metric's
useful range. Green spaces are treated as equivalent point sources at their
centroids, ignoring size, quality and public/private status; a large park reduced to
its centroid can misplace the true walking distance to its entrances, and the
sensitivity of the results to multi-entrance or area-based access points is not
tested here. OSM mapping completeness, pedestrian-network coverage, green-space classification and
topological conventions also vary across countries and are not quantified, so
data-quality heterogeneity is a genuine external-validity limit, and one the city fixed
effects do \emph{not} remove: fixed effects absorb each city's unobserved mean level,
not city-specific differences in data quality or in the morphology--resilience slope
itself. The per-city coefficients (Table~\ref{tab:hetero}) and the leave-one-city-out
transfer above are what speak to slope stability; absolute cross-city comparison of
coefficients is not warranted.
The class-exclusion sensitivity below addresses the tag-mix concern but not these
geometric and completeness ones.

\emph{Disruption and data.} Random removal is spatially uniform and does not
reproduce the clustered failures of a real flood or earthquake; targeted removal
addresses criticality but not spatial correlation. The real $100$-year flood
scenario added above makes this concrete: it materially affects only Vienna and
does not carry the distributed-disruption morphology signal, showing that resilience
to a localised hazard is a genuinely different quantity---one this paper does not
claim to explain. Earthquake and multi-hazard disruption models are left to future
work.
Targeted betweenness is approximated from $k{=}500$ sampled pivots, as exact
$O(V\lvert E\rvert)$ betweenness is intractable at city scale; the ranking is
stable in the pivot count: on the Amsterdam network, halving or
doubling $k$ (to $250$ or $1{,}000$) leaves the edge ranking essentially unchanged
(Spearman $\rho=0.985$ against $k{=}500$, with $77\%$/$90\%$ overlap of the
top-decile edges)---so the targeted scenario is not an artefact of the pivot
count. Results use a persistence-thresholded $H_0$ metric, so reported magnitudes
are denoised estimates; varying the threshold over a $4\times$ range ($0.5$--$2.0$
walk-minutes) preserves every coefficient sign and keeps the three headline
descriptors significant at $p<10^{-6}$, and the full disruption grid confirms the
signs, significances, ranking and $\rhostar$. OSM completeness varies by city and
is not quantified here, and the network-distance Vietoris--Rips variant of the
filtration is left to future work.

\emph{Modelling choices.} Two inputs could in principle shape the results: the
inclusive green-space definition and the constant $4.8$~km/h walking speed. Both were tested. The strict, accessibility-conservative green-space subset leaves the headline
signs and significance unchanged (Section~\ref{sec:results}); and, because speed
enters only as a global rescaling of the field, the walking-speed choice cannot alter
the coefficient signs, significances or the city ranking, which is confirmed
empirically ($\rho>0.98$ between the $4.8$ and $4.0$/$5.0$~km/h district AUCs). These
are therefore not load-bearing assumptions.

Finally, ``resilience'' here is a \emph{structural} property---the robustness of the
access topology to synthetic edge removal, measured on a single cross-sectional
snapshot per city---not a temporal recovery process; how access
rebuilds over time after a real disruption, and the term should be read in that
narrower, topological sense.

Read together, these caveats narrow the claim: within these eight cities, and under
the tested \emph{synthetic, distributed} edge-removal disruption (random and targeted),
street-network morphology is a consistent descriptive correlate of green-space access
resilience---robust in sign across disruption models, district resolutions, spatial
specifications and the added controls---rather than a causally identified or
cross-city-generalisable driver. Crucially, this scope is \emph{empirical}, not merely
cautionary: under the real, spatially \emph{clustered} $100$-year flood scenario the
pooled morphology effects do not reproduce (Section~\ref{sec:results}). Distributed
removal thins links roughly evenly, so the descriptors that govern how a district's
served $H_0$ structure reorganises under diffuse loss are what the metric
responds to; a clustered flood instead removes a spatially contiguous block of the
network at once, a regime in which local morphology is dominated by \emph{where} the
hazard falls rather than by routing form. The general phrasing ``street morphology
predicts urban resilience'' therefore overreaches on two counts---``urban resilience''
(the metric measures connected-access resilience, not global urban resilience) and
``predicts'' (the relationship is associational); the defensible statement is that
street morphology is associated with connected-access resilience \emph{under the
tested synthetic, distributed disruption scenarios}, and its transfer to real
clustered hazards is not established.

\section{Conclusion}
This study defines a topological metric for the resilience of green-space access, the
Wasserstein degradation curve of persistent-homology diagrams under network
disruption, and showed across $3{,}999$ districts in eight contrasting cities on four
continents that, under synthetic distributed disruption, street-network morphology is
associated with how robustly a city's
pedestrian network keeps its \emph{connected} served locations close to nature. Two
effects are robust throughout---longer street segments (the largest) and more
circuitous routing are associated with greater resilience, and both survive a wild
cluster bootstrap and spatial-error and spatial-lag models---while higher orientation
entropy is associated with less resilience but proves the least dependable, weakening
across cities and vanishing once diagram magnitude is controlled. These associations
hold for distributed disruption and do not carry to the real clustered-flood scenario. The contribution has two boundaries. The relationship
is an \emph{association}, not a causal effect: the design is cross-sectional and
observational. And it is an association with the \emph{proposed} topological
indicator; whether that indicator is a fully valid measure of resilience---rather
than a well-behaved, configuration-sensitive proxy---is the study's central open
question. Eight cities also still furnish only eight clusters (see
Limitations). The paper is precise
about what the metric captures: the graded topological persistence of the
\emph{connected} green-access structure under progressive network thinning---moderately
and heterogeneously validated within city against a distance-based accessibility
benchmark across all eight cities, and shown by its negative correlation with
connectivity degradation to be a distinct construct from percolation robustness rather
than total accessibility including outright disconnection. The mathematics of persistence turns a
static SDG~11.7 indicator into a dynamic, form-aware one. Future work will
extend the real-hazard analysis beyond river floods to earthquakes and multi-hazard
scenarios, and add a network-metric Vietoris--Rips filtration, toward a
resilience-aware accessibility standard for
sustainable cities.

\paragraph{Reproducibility.} All data are open (OpenStreetMap, GHS-UCDB) and the
pipeline is open-source Python (OSMnx, NetworkX, GeoPandas, GUDHI/\texttt{persim},
H3, \texttt{statsmodels}); every reported number is recorded in the project's
results log and traces to a committed artifact.


\appendix
\section{Robustness and sensitivity analyses}\label{app:robust}
This appendix collects the diagnostics summarised in the Results. Throughout, the
three headline descriptors are circuity, mean street length and orientation entropy.

\paragraph{Spatial dependence.} Districts nest within cities and are spatially
contiguous, so the fixed-effects OLS residuals are not independent. Global Moran's~$I$
on first-ring H3 adjacency is positive and significant in every city (up to $0.18$ in
Barcelona and Singapore, mostly permutation $p\le0.01$), and Anselin
Lagrange-multiplier tests on the OLS residuals
decisively reject the non-spatial model (LM-error $p<10^{-135}$, LM-lag
$p<10^{-177}$; the robust LM statistics favour the lag form, $227$ vs.\ $31$), with
both spatial models cutting the AIC by more than $500$ relative to OLS
($22{,}063\to21{,}475$ for the error model, $21{,}376$ for the lag model), so
spatial-error and spatial-lag models are estimated spatial-error and spatial-lag models (Table~\ref{tab:spatial}).
Both absorb strong dependence ($\hat\lambda=0.53$, $\hat\rho=0.51$), yet the three
headline descriptors keep their sign and stay strongly significant ($|z|\ge5$).
City-clustered standard errors (Table~\ref{tab:spatial}), the most conservative
treatment of the eight-cluster design, leave all three significant, including
orientation entropy.

\begin{table}[t]\centering
\caption{Inference robustness of the eight-city district regression: city-clustered
standard errors (eight clusters, the most conservative treatment of the nesting) and
estimated spatial-error and spatial-lag models (ML, within-city H3 weights).
Coefficients are on AUC. All three headline descriptors survive city clustering and
both spatial models, despite substantial estimated spatial dependence.}
\label{tab:spatial}
\begin{tabular}{lrrrr}
\toprule
Morphology descriptor & coef (OLS) & $p$ (cluster) & Sp.\ error $z$ & Sp.\ lag $z$\\
\midrule
circuity & $-7.39$ & $2.2\times10^{-4}$ $^{***}$ & $-5.2$ & $-6.6$\\
mean street length & $-0.050$ & $4.5\times10^{-19}$ $^{***}$ & $-7.5$ & $-8.8$\\
orientation entropy & $+1.85$ & $0.036$ $^{*}$ & $+10.1$ & $+9.6$\\
intersection density & $+0.02$ & $0.89$ & $-1.5$ & $-1.2$\\
\midrule
spatial parameter & & & $\hat\lambda=0.53$ & $\hat\rho=0.51$\\
\bottomrule
\end{tabular}
\end{table}

\paragraph{Multicollinearity and cross-city heterogeneity.} The headline descriptors
are not collinear: within-city, the largest pairwise correlation among circuity, mean
street length and orientation entropy is $0.31$, and all variance-inflation factors
are below three (the only strong pair is mean street length and intersection density,
$r=-0.80$, which is why intersection density is unstable and not relied on). The
pooled slopes are also not a pooling artefact: fitting the regression separately per
city, each headline sign holds in seven of the eight cities (Table~\ref{tab:hetero}),
with Nairobi the exception for circuity and Phoenix for orientation entropy.
The magnitudes, however, are \emph{not} constant across cities, and this is stated
explicitly rather than assumed away. A joint test of the morphology$\times$city
interactions rejects slope homogeneity ($F=12.5$ on $21$ terms,
$p<10^{-41}$), and the per-city slopes are widely dispersed (for circuity and mean
street length the between-city standard deviation is of the same order as the mean
slope; coefficient of variation $\approx1.0$). A random-slopes mixed model
(standardised) agrees: the mean slope and its between-city SD are
$-0.31\pm0.17$ for mean street length, $+0.28\pm0.21$ for orientation entropy and
$-0.14\pm0.13$ for circuity---directionally consistent but with real magnitude
heterogeneity. City fixed effects remove only mean-level differences, not this
variation in slope. The pooled coefficient is therefore interpreted throughout as the
\emph{average within-city conditional association across these eight cities}, not as a
homogeneous or universal effect of morphology on resilience; mean street length is the
most stable (largest mean-to-SD ratio), circuity the least.

\begin{table}[t]\centering
\caption{Per-city regression coefficients (AUC on morphology, one city at a time).
Each headline sign---circuity ($-$), mean street length ($-$), orientation entropy
($+$)---holds in seven of the eight cities.}
\label{tab:hetero}
\begin{tabular}{lrrr}
\toprule
City & circuity & mean street length & orientation entropy\\
\midrule
Amsterdam & $-4.07$ & $-0.04$ & $+1.17$\\
Barcelona & $-6.57$ & $+0.02$ & $+6.17$\\
Bogotá    & $-18.62$ & $-0.21$ & $+1.46$\\
İstanbul  & $-14.51$ & $-0.03$ & $+3.59$\\
Nairobi   & $+4.86$ & $-0.05$ & $+2.56$\\
Phoenix   & $-3.70$ & $-0.06$ & $-0.72$\\
Singapore & $-4.12$ & $-0.04$ & $+4.58$\\
Vienna    & $-5.56$ & $-0.04$ & $+3.39$\\
\bottomrule
\end{tabular}
\end{table}

\paragraph{Out-of-sample cross-city transfer (leave-one-city-out).} With only eight
clusters, the sharper question is whether the relationship \emph{transfers} to a city
the model never saw. A leave-one-city-out test addresses this directly: within-city
slopes are fitted on seven cities and used to predict the ordering of districts in the
held-out eighth (the held-out city's own mean is removed, so the target is its
within-city ranking, which needs no estimated fixed effect). The transfer is positive
in \emph{all eight} held-out cities, with mean out-of-sample Spearman $+0.42$
(Barcelona $+0.61$, Singapore $+0.58$, Vienna $+0.49$, İstanbul $+0.47$, Nairobi
$+0.39$, Amsterdam $+0.36$, Bogotá $+0.34$, and Phoenix the weak case at $+0.07$,
$p=0.02$). Morphology fitted elsewhere therefore predicts, imperfectly but consistently,
the resilience ranking of an unseen city's districts---the strongest external-validity
evidence the eight-city design can offer, short of an independent validation sample.

\paragraph{Disruption scenarios (five-city core).} Under \emph{targeted} disruption (removal of
high-betweenness edges, ranked by $k{=}500$-sampled betweenness) every coefficient
sign in Table~\ref{tab:regression} is reproduced: circuity $-5.77$
($p=4.6\times10^{-10}$), mean street length $-0.034$ ($p=2.7\times10^{-13}$),
orientation entropy $+0.50$ ($p=1.6\times10^{-3}$), and intersection density
(not significant under the full random grid) reaching $p=0.022$ ($-0.135$). Absolute
AUC values are lower (Amsterdam mean AUC $2.03$ vs.\ $4.46$ under random) but the
ranking is unchanged.

\paragraph{Real flood-hazard scenario.} The earlier low-elevation
proxy with a \emph{real} $100$-year river-flood scenario, using the open JRC Global
Flood Hazard maps \citep{dottori2016globalflood} (floodMapGL, water depth per
$\approx1$~km cell): a node is flood-exposed if it lies in a cell with $\ge0.5$~m of
inundation, and the hazard removal progressively cuts flood-exposed edges. Two
honest findings follow. First, real $100$-year exposure among these urban centres is
strongly city-specific: it is material only in Vienna ($23\%$ of nodes, on the
Danube floodplain) and negligible elsewhere (Phoenix $0.7\%$, Amsterdam $0.4\%$,
Barcelona $0.3\%$; İstanbul, Bogotá, Singapore and Nairobi essentially $0\%$, their
centres sitting off the modelled river-flood plain). Second, where the scenario
bites (Vienna), the morphology signal established for \emph{distributed} disruption
does \emph{not} cleanly carry over to this \emph{spatially clustered} flood: in a
pooled flood-scenario regression the morphology coefficients are non-significant,
and within Vienna only circuity keeps a weak same-signed effect (Spearman $-0.13$,
$p=0.006$) while orientation entropy even reverses. This is consistent with, and
concretely demonstrates, the caveat that resilience to network-wide disruption is a
different quantity from resilience to a localised, geographically concentrated
hazard; morphology governs the former, not necessarily the latter.

\paragraph{Disruption-grid and resolution sensitivity (five-city core).} A coarser $5\times3$
disruption grid leaves every sign unchanged and keeps the three headline descriptors
significant at $p<10^{-3}$ (circuity $-8.99$, mean street length $-0.0478$,
orientation entropy $+1.17$), with the city AUC ranking and $\rhostar$ values within
$\approx0.02$ of the full $8\times10$ grid. The one descriptor that moves is
intersection density: significant on the coarse grid ($-0.199$,
$p=8.6\times10^{-3}$) but not on the full grid ($p=0.22$). Re-tiling at H3
resolutions $7$ ($485$ districts) and $9$ ($13{,}782$ districts) likewise leaves the
three headline descriptors significant at $p<10^{-3}$ with constant signs across a
$28\times$ range of district counts; magnitudes scale with cell size but no sign
flips. Figure~\ref{fig:robustness} summarises both axes.

\begin{figure}[t]\centering
\fbox{\parbox[c][4.5cm][c]{0.72\linewidth}{\centering\ttfamily fig\_robustness\\[2pt]\normalfont\itshape (drag-drop this PDF into figures/, then restore \textbackslash includegraphics)}}
\caption{Robustness of the morphology$\leftrightarrow$resilience relationship.
\textbf{(a)} Regression coefficients (with $95\%$ confidence intervals) for each
morphology descriptor under the two distributed disruption scenarios (random and
targeted) at H3 resolution~8; the coefficient signs agree, and the three headline
descriptors (circuity, mean street length, orientation entropy) stay significant,
while intersection density crosses zero under random disruption. \textbf{(b)} Sign and
significance of each coefficient across H3 resolutions $7$/$8$/$9$: the three
headline descriptors stay significant with a constant sign, while intersection
density, the weakest and grid-sensitive effect, is not significant at the headline
full-grid specification.}
\label{fig:robustness}
\end{figure}

\paragraph{Persistence-threshold sensitivity (five-city core).} Varying the $H_0$ denoising
threshold over a $4\times$ range leaves the headline unchanged: at $0.5$, $1.0$ and
$2.0$~walk-minutes the three descriptors keep their sign and stay significant at
$p<10^{-6}$ (circuity $-11.20/-10.16/-8.03$; mean street length
$-0.050/-0.048/-0.042$; orientation entropy $+1.10/+1.01/+0.93$). Magnitudes shrink
smoothly as more denoising is applied but never change sign, and intersection
density stays weak (ns at $0.5$ and $1.0$, $p=0.04$ at $2.0$). The betweenness pivot
count is equally inconsequential: on the Amsterdam network, halving or doubling $k$
(to $250$ or $1{,}000$) leaves the edge ranking essentially unchanged (Spearman
$\rho=0.985$ against $k{=}500$, with $77\%$/$90\%$ overlap of the top-decile edges).

\paragraph{Omitted-variable controls (five-city core).} City fixed effects absorb only between-city
confounds, so the district-level controls that the data support are added---total
green-space area, the share of edges on the trunk/primary/secondary/tertiary road
hierarchy, within-district relief (elevation standard deviation), and distance to
the urban-centre centroid---and refit (Table~\ref{tab:controls}). Three of the four
controls are themselves significant (more green area and greater centre distance
with greater resilience, more major-road edges with less), yet the three headline
descriptors keep their sign and stay significant at $p<10^{-11}$ (circuity $-9.55$,
mean street length $-0.039$, orientation entropy $+1.11$). Population and built
density, land-use mix, socioeconomic status and pedestrian-infrastructure quality
require external data and remain uncontrolled (Limitations).

\begin{table}[t]\centering
\caption{Omitted-variable robustness (five-city core): district regression before and
after adding four controls, $n=2{,}603$, with city fixed effects. The three headline
morphology effects keep their sign and significance; the added controls explain
additional variance ($R^2$ $0.17\to0.21$). Sig.: $^{***}p<0.001$, $^{**}p<0.01$.}
\label{tab:controls}
\begin{tabular}{lrr}
\toprule
& base & $+$controls\\
Variable & coef ($p$) & coef ($p$)\\
\midrule
circuity & $-10.16$ $^{***}$ & $-9.55$ $^{***}$\\
mean street length & $-0.048$ $^{***}$ & $-0.039$ $^{***}$\\
orientation entropy & $+1.01$ $^{***}$ & $+1.11$ $^{***}$\\
intersection density & $-0.09$ & $-0.10$\\
\midrule
green-space area (km$^2$) & --- & $-2.25$ $^{**}$\\
major-road share & --- & $+4.21$ $^{***}$\\
relief (m) & --- & $-0.008$\\
centre distance (km) & --- & $-0.107$ $^{***}$\\
\bottomrule
\end{tabular}
\end{table}


\section*{Declarations}

\paragraph{Funding.} The author received no specific funding for this work.

\paragraph{Competing interests.} The author declares no competing interests.

\paragraph{Data availability.} All input data are open and publicly available.
Pedestrian networks and green spaces were obtained from OpenStreetMap via
OSMnx; urban-centre boundaries from the GHS Urban Centre Database (GHS-UCDB
R2019A, European Commission Joint Research Centre); residential population from the
GHS-POP grid (R2023A, $1$~km) for the population-weighted validation; $100$-year
river-flood depth from the open JRC Global Flood Hazard maps (floodMapGL) for the
real-flood scenario; and elevation data from the public AWS Terrain Tiles service. No new data were collected. The derived
artifacts underlying every reported figure and table (per-district resilience and
morphology tables, regression outputs, and city-level resilience curves) are
recorded in the project's results log, released with the code at
[repository URL omitted for double-blind peer review].

\paragraph{Code availability.} The complete analysis pipeline is implemented in
open-source Python (OSMnx, NetworkX, GeoPandas, GUDHI, \texttt{persim}, H3,
\texttt{statsmodels}) and is publicly available at
[repository URL omitted for double-blind peer review], together with the
configuration files and the results log that reproduce every reported number.

\paragraph{Author contributions.} The author conceived
the study, developed the method and software, performed the analysis, and wrote
the manuscript.

\bibliography{sn-bibliography}

\end{document}